\documentclass[11pt]{article}

\usepackage[margin=1in]{geometry}
\usepackage{amsmath,amssymb,amsthm,bm,booktabs,array}
\usepackage[colorlinks=true,linkcolor=blue,urlcolor=blue]{hyperref}
\usepackage{microtype}

\newcommand{\SCA}{\mathsf{SCA}}
\newcommand{\Ff}{\mathsf F}
\newcommand{\Vir}{\mathsf{Vir}}
\newcommand{\one}{\mathbf 1}
\newcommand{\cF}{\mathcal F}
\newcommand{\cV}{\mathcal V}
\newcommand{\hatrho}{\widehat\rho}
\newcommand{\se}{s_{\rm e}}
\newcommand{\so}{s_{\rm o}}
\newcommand{\bpz}[2]{\bigl\langle\!\bigl\langle #1\bigm|#2\bigr\rangle\!\bigr\rangle}

\title{The NS Sphere Block from Free Fermion plus Super-Virasoro\\
and from Two Virasoro Algebras}
\author{Machine note}
\date{20 August 2026}

\begin{document}
\maketitle

\begin{abstract}
We derive the two-Virasoro expansion of the sphere four-point block of four
bottom NS superconformal primaries, using exactly the slot order and odd
three-form sign of \texttt{Human Notes/SCblock.tex}.  The auxiliary fermion is
put in the identity state at every external puncture.  With the ungraded
pairing on \(\Ff_{1/2}\otimes\SCA_{1/2}\), its block is one and the enlarged
block equals the super-Virasoro block.  Resolving the same block in the
\(\Vir\oplus\Vir\) branching basis gives a sum over
\(n\in\frac12\mathbb Z\).  The human-note endpoint convention supplies the
essential factor \((-1)^{2n}\).  A direct PBW computation agrees with the resulting formula
through even level two and odd level \(3/2\).  The nonchiral crossing equation
is unchanged because this sign is common to the whole odd chiral block and
therefore squares to one.
\end{abstract}

\section{Conventions and the ordered NS block}

We use the ordinary central charge and momentum convention of the human note,
\begin{equation}
 Q=b+b^{-1},\qquad
 c=\frac32+3Q^2,\qquad
 h(P)=\frac12\left(\frac{Q^2}{4}-P^2\right).
 \label{eq:sca-parameters}
\end{equation}
The Liouville continuum is reached by \(P=i p\); it is useful to keep \(P\)
generic while proving the algebraic identity.  All external states
\(\phi_i\) are even bottom components.  The ordered correlator is
\begin{equation}
 \left\langle
 \phi_4(\infty)\phi_3(1)\phi_2(z)\phi_1(0)
 \right\rangle .
 \label{eq:ordered-correlator}
\end{equation}

Let \(\xi_A=\mathbf L_{-A}\phi_h\) be an ordered NS PBW state, let
\(|A|\) be its level, and let \(p(A)\) be its fermion parity.  The inverse
super-Gram matrix is \(\mathbf G_h^{AB}\).  We strip the common leading
power and define the two parity blocks by
\begin{align}
 F_a^{4321}(h;z)
 &=z^{h-h_1-h_2}\cF_a^{4321}(h;z),
 \nonumber\\
 \cF_a^{4321}(h;z)
 &=\sum_{\substack{|A|=|B|\\p(A)=p(B)=a}}
 z^{|A|}\,
 \rho_a(\phi_4,\phi_3,\xi_A)\,
 \mathbf G_h^{AB}\,
 \rho_a(\xi_B,\phi_2,\phi_1),
 \qquad a\in\{0,1\}.
 \label{eq:sca-block-definition}
\end{align}
Thus \(\cF_0\) has integer powers and \(\cF_1\) has half-odd-integer
powers.

The sign which matters below is fixed by the component table in the human
note.  For even primaries,
\begin{equation}
 \rho_1(G_{-1/2}\phi_1,\phi_2,\phi_3)=+1,
 \qquad
 \rho_1(\phi_1,G_{-1/2}\phi_2,\phi_3)=+1,
 \qquad
 \boxed{\rho_1(\phi_1,\phi_2,G_{-1/2}\phi_3)=-1.}
 \label{eq:human-odd-table}
\end{equation}
The last minus sign is retained throughout this note.  Some production
sphere-block code uses \(+1\) in the third entry; that is a different
chiral convention and differs from the present odd block by an overall
minus sign.

\section{The ungraded tensor-product pairing}

Adjoin an auxiliary NS Majorana fermion with
\begin{equation}
 \{\psi_r,\psi_s\}=\delta_{r+s,0},
 \qquad \psi_r^\dagger=-\psi_{-r},
 \qquad r\in\mathbb Z+\frac12.
 \label{eq:fermion-algebra}
\end{equation}
The internal enlarged module is
\begin{equation}
 \mathcal M_{\rm NS}(h)
 =\mathcal V_{\Ff_{1/2}}\otimes\mathcal V_{\SCA_{1/2}}(h).
 \label{eq:enlarged-module}
\end{equation}
The essential choice is the prime, ungraded pairing
\begin{equation}
 \boxed{
 \bpz{u\otimes x}{v\otimes y}'
 =\bpz{u}{v}_{\Ff}\,\bpz{x}{y}_{\SCA}.}
 \label{eq:ungraded-pairing}
\end{equation}
There is no Koszul sign crossing the two tensor factors.  This differs from
the standard graded tensor-product pairing.

For product states the normalized three-form is the one printed in the
current human note,
\begin{equation}
 \hatrho_a(u_1\otimes x_1,u_2\otimes x_2,u_3\otimes x_3)
 =(-1)^{|x_1|(|u_2|+|u_3|)+|x_2||u_3|}
 \rho^{\Ff}(u_1,u_2,u_3)\rho_a(x_1,x_2,x_3).
 \label{eq:three-form-factorization}
\end{equation}

At the external punctures choose
\(\widehat\phi_i=\one_{\Ff}\otimes\phi_i\).  In the product PBW basis the
Gram matrix and its inverse factor as ordinary tensor products because of
Eq.~\eqref{eq:ungraded-pairing}.  Equation~\eqref{eq:three-form-factorization}
has no crossing sign at either sphere trinion: the external SCA states and
the external fermion vacua are even.  Moreover,
\begin{equation}
 \rho^{\Ff}(\one,\one,u)=\rho^{\Ff}(u,\one,\one)=0
 \quad\text{unless}\quad u=\one.
 \label{eq:fermion-identity-selection}
\end{equation}
It follows coefficient by coefficient that
\begin{equation}
 \widehat{\cF}_a^{4321}(h;z)
 =\cF_{\Ff}^{1111}(z)\,\cF_a^{4321}(h;z),
 \qquad
 \boxed{\cF_{\Ff}^{1111}(z)=1,\quad
 \widehat{\cF}_a^{4321}=\cF_a^{4321}.}
 \label{eq:sphere-factorization}
\end{equation}
Unlike the genus-two theta graph, there is no multivariable polarized
convolution: the free-fermion identity block kills every nonvacuum
auxiliary state.

\section{The same enlarged block in the double-Virasoro basis}

For \(b\ne1\), the human-note generators give
\(\Vir^{(1)}\oplus\Vir^{(2)}\subset\Ff\oplus\SCA\).  Their central charges
are
\begin{equation}
 c^{(i)}=1+6\bigl(Q^{(i)}\bigr)^2,
 \qquad
 (b^{(1)})^2=\frac{2b^2}{1-b^2},
 \qquad
 (b^{(2)})^{-2}=\frac{2}{b^2-1},
 \qquad Q^{(i)}=b^{(i)}+(b^{(i)})^{-1}.
 \label{eq:double-central-charges}
\end{equation}
In the notation of the human note, the NS branching rule is
\begin{equation}
 \mathcal M_{\rm NS}(h(P))
 \cong\bigoplus_{n\in\frac12\mathbb Z}
 \mathcal V_{c^{(1)}}(h_n^{(1)}(P))\otimes
 \mathcal V_{c^{(2)}}(h_n^{(2)}(P)),
 \qquad
 h_n^{(1)}+h_n^{(2)}=h(P)+2n^2.
 \label{eq:branching-rule}
\end{equation}
In a square-root-free parametrization useful for computation,
\begin{align}
 h_n^{(1)}(P)
 &=\frac{(Q^{(1)})^2}{4}
 -\frac{(P+2nb)^2}{2-2b^2},
 &
 h_n^{(2)}(P)
 &=\frac{(Q^{(2)})^2}{4}
 -\frac{(P+2n/b)^2}{2-2b^{-2}}.
 \label{eq:branch-weights}
\end{align}
The branch primary \(v_n(P)\) has parity \(2n\pmod2\) and ungraded norm
\(N_n(P)=\bpz{v_n}{v_n}'\).

Now the reason for Eq.~\eqref{eq:ungraded-pairing} becomes explicit.  If
\(I,J\) and \(I',J'\) are ordinary Virasoro partitions, then
\begin{align}
 &\bpz{L_{-I}^{(1)}L_{-J}^{(2)}v_n}
 {L_{-I'}^{(1)}L_{-J'}^{(2)}v_n}'
 \nonumber\\
 &\hspace{2cm}=
 N_n(P)\,(G_{h_n^{(1)}}^{(1)})_{II'}
 (G_{h_n^{(2)}}^{(2)})_{JJ'}.
 \label{eq:double-gram-factorization}
\end{align}
The inverse therefore factorizes with no parity operator left on the tube.
The graded product pairing would insert a state-dependent crossing sign and
would not give Eq.~\eqref{eq:double-gram-factorization} in this basis.

Let \(\cV(c;d_4,d_3,d_2,d_1;d;z)=1+O(z)\) denote the reduced ordinary
Virasoro sphere block in the same slot order.  Define the two unnormalized
branching vertices
\begin{equation}
 T_{43}(n)=\hatrho_{2n\bmod2}(v_{4,0},v_{3,0},v_n),
 \qquad
 T_{21}(n)=\hatrho_{2n\bmod2}(v_n,v_{2,0},v_{1,0}).
 \label{eq:endpoint-vertices}
\end{equation}
We now use exactly the branching-coefficient notation of the human note:
\begin{equation}
 \boxed{
 B_a(h_1,h_2,h_3;n_1,n_2,n_3)
 =\frac{
 \hatrho_a(v_{1,n_1},v_{2,n_2},v_{3,n_3})
 }{
 \lVert v_{1,n_1}\rVert
 \lVert v_{2,n_2}\rVert
 \lVert v_{3,n_3}\rVert
 }.}
 \label{eq:human-B-definition}
\end{equation}
The external branch vectors have unit norm.  Hence the two coefficients at
the sphere endpoints are
\begin{equation}
 B_a(h_4,h_3,h;0,0,n)=\frac{T_{43}(n)}{\lVert v_n\rVert},
 \qquad
 B_a(h,h_2,h_1;n,0,0)=\frac{T_{21}(n)}{\lVert v_n\rVert}.
 \label{eq:sphere-B-endpoints}
\end{equation}
Ordinary Virasoro Ward identities act independently in the two factors.
Inserting the inverse of Eq.~\eqref{eq:double-gram-factorization} into the
four-point block gives
\begin{align}
 \boxed{
 \cF_a^{4321}(h(P);z)
 =\sum_{\substack{n\in\frac12\mathbb Z\\(-1)^{2n}=(-1)^a}}
 z^{2n^2}
 B_a(h_4,h_3,h;0,0,n)
 B_a(h,h_2,h_1;n,0,0)
 \cV_n^{(1)}(z)\cV_n^{(2)}(z).}
 \label{eq:final-B-formula}
\end{align}
where
\begin{equation}
 \cV_n^{(i)}(z)=
 \cV\!\left(
 c^{(i)};
 h_{0}^{(i)}(P_4),h_{0}^{(i)}(P_3),
 h_{0}^{(i)}(P_2),h_{0}^{(i)}(P_1);
 h_n^{(i)}(P);z\right).
 \label{eq:factor-blocks}
\end{equation}
The product of branching coefficients is independent of a rephasing of
\(v_n\).  If the same square-root branch for \(\lVert v_n\rVert\) is used at
both ends, it can equivalently be written
\begin{equation}
 B_a(h_4,h_3,h;0,0,n)
 B_a(h,h_2,h_1;n,0,0)
 =\frac{T_{43}(n)T_{21}(n)}{N_n(P)}.
 \label{eq:B-product-phase-independent}
\end{equation}
Thus \(T_{43}T_{21}/N_n\) is only a phase-independent representation of the
human-note product \(B_aB_a\), not a replacement notation for the branching
coefficient.

\section{The human-note endpoint sign and the explicit blow-up formula}

We now translate the branching coefficients in
Eq.~\eqref{eq:final-B-formula} to the factorized product \(l\) used in the
human note and in the blow-up literature.  Only in the following product
formula, \(m,k',k\in\mathbb Z\) are integer Fermi-sea labels; the internal
sphere label is \(k=2n\).  For \(\sigma,\tau=\pm1\), set
\begin{equation}
 x_{\sigma\tau}=\alpha+\sigma P'+\tau P,
 \qquad
 r_{\sigma\tau}=\frac{m+\sigma k'+\tau k}{2}.
 \label{eq:x-r}
\end{equation}
For \(r\ge0\), define
\begin{align}
 \se(x,r)
 &=2^{-r^2/2}
 \prod_{\substack{i,j\ge1,\ i+j\le2r\\i+j\equiv0\ (2)}}
 [x+(i-1)b+(j-1)b^{-1}],
 \nonumber\\
 \so(x,r)
 &=2^{-r(r+1)/2}
 \prod_{\substack{i,j\ge1,\ i+j\le2r+1\\i+j\equiv1\ (2)}}
 [x+(i-1)b+(j-1)b^{-1}],
 \label{eq:s-functions}
\end{align}
and continue to negative \(r\) by
\begin{equation}
 \se(x,r)=(-1)^r\se(Q-x,-r),
 \qquad
 \so(x,r)=\so(Q-x,-r).
 \label{eq:s-continuation}
\end{equation}
With \(\operatorname{Int}(x)=\operatorname{sgn}(x)\lfloor|x|\rfloor\),
\begin{equation}
 l(\alpha,m\mid P',k',P,k)
 =\begin{cases}
 \displaystyle\prod_{\sigma,\tau=\pm1}
 \se(x_{\sigma\tau},r_{\sigma\tau}),
 &m+k+k'\text{ even},\\[0.8em]
 \displaystyle\prod_{\sigma,\tau=\pm1}
 \so(x_{\sigma\tau},\operatorname{Int}r_{\sigma\tau}),
 &m+k+k'\text{ odd}.
 \end{cases}
 \label{eq:l-product}
\end{equation}
For the branch label \(n\in\frac12\mathbb Z\), the norm is the identity
specialization with the integer Fermi-sea label \(2n\):
\begin{equation}
 N_n(P)=\se(2P,2n)\se(-2P,-2n).
 \label{eq:branch-norm}
\end{equation}

For the sphere endpoints define
\begin{align}
 L_{43}(n)
 &=l\left(\frac Q2+P_3,0\mid P_4,0,P,2n\right),
 \nonumber\\
 L_{21}(n)
 &=l\left(\frac Q2+P_2,0\mid P,2n,P_1,0\right).
 \label{eq:sphere-l-factors}
\end{align}
The human-note odd table \eqref{eq:human-odd-table} fixes the ordered endpoint
translation:
\begin{equation}
 \boxed{T_{43}(n)=(-1)^{2n}L_{43}(n),\qquad
 T_{21}(n)=L_{21}(n).}
 \label{eq:endpoint-translation}
\end{equation}
This is not an optional phase of \(v_n\): a rephasing of \(v_n\) occurs at
both endpoints and cancels against \(N_n\).  It is the relative sign between
placing the odd fixed-parity tensor in slot three and in slot one.  Direct
branching-vector calculations verify Eq.~\eqref{eq:endpoint-translation} for
\(n=0,\pm\frac12\), and for the one-active-leg cases \(n=1,\frac32\); the
general factor is fixed by the same parity-homogeneous three-form
normalization.

Equations~\eqref{eq:sphere-B-endpoints} and
\eqref{eq:endpoint-translation} give the explicit product of the two
human-note branching coefficients:
\begin{equation}
 \boxed{
 B_a(h_4,h_3,h;0,0,n)B_a(h,h_2,h_1;n,0,0)
 =(-1)^{2n}\frac{L_{43}(n)L_{21}(n)}{N_n(P)}.}
 \label{eq:explicit-B-product}
\end{equation}
Substituting Eq.~\eqref{eq:explicit-B-product} in the boxed block relation
\eqref{eq:final-B-formula} reproduces the square-root-free expression used in
the numerical check.  Since \((-1)^{2n}=(-1)^a\) within a fixed block, the
sign is a common minus on the odd chiral block and is one on the even block.
Equation~\eqref{eq:final-B-formula} is first derived for generic \(b\ne1\)
and generic momenta; the \(b=1\) result is its analytic limit after the two
Virasoro factors have been combined.

\subsection{Analytic \texorpdfstring{level-\(1/2\)}{level-1/2} sign check}

The first branch primaries are
\begin{equation}
 v_{\pm\frac12}(P)=
 \left[G_{-1/2}+\left(\frac Q2\pm P\right)\psi_{-1/2}\right]\phi,
 \label{eq:first-branch-vectors}
\end{equation}
and their ungraded norms are
\begin{equation}
 N_{+\frac12}(P)=-P(Q+2P),
 \qquad
 N_{-\frac12}(P)=P(Q-2P).
 \label{eq:first-branch-norms}
\end{equation}
The external fermion identities kill the \(\psi_{-1/2}\) term.  Hence
Eq.~\eqref{eq:human-odd-table} gives \(T_{43}(\pm\frac12)=-1\),
\(T_{21}(\pm\frac12)=+1\), while
\(L_{43}(\pm\frac12)=L_{21}(\pm\frac12)=1\).  The double-Virasoro answer at
the first odd level is therefore
\begin{align}
 [z^{1/2}]\cF_1
 &=-\left(\frac1{N_{+\frac12}}+\frac1{N_{-\frac12}}\right)
 =-\frac{4}{Q^2-4P^2}
 =\boxed{-\frac1{2h(P)}}.
 \label{eq:level-half-check}
\end{align}
Direct super-Virasoro sewing gives the same result:
\((-1)(+1))/\lVert G_{-1/2}\phi\rVert^2=-1/(2h)\).
Dropping \((-1)^{2n}\) gives \(+1/(2h)\), the opposite convention.

\section{Direct low-order PBW check}

The reproducible calculation is
\begin{center}
 \path{Code/check_ns_sphere_double_virasoro.py}.
\end{center}
It constructs the finite-\(c\) NS PBW basis and Gram matrix directly,
evaluates both descendant three-forms from Ward identities, and compares
with Eq.~\eqref{eq:final-B-formula}.  Ordinary Virasoro Gram matrices
are evaluated independently through level two.  The generic test point is
\begin{gather}
 b=1.37,\qquad Q=2.09992700729927018,
 \nonumber\\
 c=14.7290803079546091,\qquad P=0.47,
 \nonumber\\
 (P_1,P_2,P_3,P_4)=(0.23,-0.31,0.41,-0.19),\qquad
 h(P)=0.440761679498108683.
 \label{eq:numeric-point}
\end{gather}
The reduced coefficients are
\begin{center}
\begin{tabular}{@{}rrrrr@{}}
\toprule
\(2L\)&direct human sign&double Virasoro&absolute error&without \((-1)^{2n}\)\\
\midrule
0&\( 1.000000000000000\)&\( 1.000000000000000\)&0&\( 1.000000000000000\)\\
1&\(-1.134399888323653\)&\(-1.134399888323653\)&\(5.04\times10^{-17}\)&\( 1.134399888323653\)\\
2&\( 0.178198040229849\)&\( 0.178198040229849\)&\(1.54\times10^{-17}\)&\( 0.178198040229849\)\\
3&\(-0.484955260912690\)&\(-0.484955260912690\)&\(1.75\times10^{-16}\)&\( 0.484955260912690\)\\
4&\( 0.120205847599768\)&\( 0.120205847599768\)&\(7.79\times10^{-17}\)&\( 0.120205847599768\)\\
\bottomrule
\end{tabular}
\end{center}
Thus the check reaches even level two, including the new \(n=\pm1\) branch
primaries, and odd level \(3/2\).  The maximum discrepancy is
\(1.75\times10^{-16}\).  Omitting the human endpoint sign flips every odd
coefficient and fails at the first possible order.

\section{Crossing equation}

For four bottom NS Liouville primaries the nonchiral direct-channel
decomposition in the project convention is
\begin{align}
 G_{4321}^{(s)}(z,\bar z)
 =\int_0^\infty\frac{dp}{\pi}\bigl[&
 C_{12p}C_{34p}\,F_0^{4321}(p;z)F_0^{4321}(p;\bar z)
 \nonumber\\[-0.2em]
 &+\widetilde C_{12p}\widetilde C_{34p}\,
 F_1^{4321}(p;z)F_1^{4321}(p;\bar z)\bigr].
 \label{eq:s-channel-crossing}
\end{align}
The crossed channel is obtained by
\((P_1,P_2,P_3,P_4;z)\mapsto(P_3,P_2,P_1,P_4;1-z)\).

Equation~\eqref{eq:final-B-formula} changes the odd chiral convention by
\(F_1^{\rm human}=-F_1^{\rm production}\), while leaving \(F_0\) unchanged.
The same sign occurs in the antiholomorphic block, so
\begin{equation}
 F_1^{\rm human}(z)F_1^{\rm human}(\bar z)
 =F_1^{\rm production}(z)F_1^{\rm production}(\bar z).
 \label{eq:odd-sign-cancels}
\end{equation}
Consequently crossing is exactly insensitive to this simultaneous chiral
sign flip.  It does, however, test the relative coefficient of the even and
odd \emph{nonchiral} families.  We therefore performed a new computation in
which the two continuum channels were evaluated independently, without
averaging, rescaling, or imposing equality.  The driver is
\begin{center}
 \path{Code/stress_ns_crossing_extended.py}.
\end{center}
It works at the self-dual point \(b=1\), where the physical structure
constants used here are implemented.  Changing only the block central charge
would not define a check at another \(b\), so no such artificial \(b\)-scan
was made.

The blocks use the direct functional \(c\)-recursion at \(c=27/2\), with
exact global \(\mathfrak{osp}(1|2)\) leaves and no local-\(z\) or elliptic-
\(q\) truncation.  We used 80-digit block arithmetic, 35-digit structure
constants, \(p\le5\), and 32-point Gauss--Legendre quadrature.  For each of
the three momentum quadruples
\begin{equation}
 \left(\frac12,\frac13,\frac14,\frac35\right),\qquad
 (0.20,0.45,0.70,0.35),\qquad
 (0.80,0.15,0.55,0.40),
 \label{eq:crossing-momenta}
\end{equation}
we evaluated seven real points \(z=0.2,0.3,\ldots,0.8\) and three complex
points \(z=0.3+0.2i,0.5+0.25i,0.7+0.2i\).  Thus the main ledger contains
90 raw channel comparisons.  With
\begin{equation}
 R=\frac{|G^{(s)}(z,\bar z)-G^{(t)}(1-z,1-\bar z)|}
 {\max(|G^{(s)}(z,\bar z)|,|G^{(t)}(1-z,1-\bar z)|)},
 \label{eq:crossing-residual}
\end{equation}
the maximum over the ten points is
\begin{center}
\begin{tabular}{@{}lrrrr@{}}
\toprule
momenta in Eq.~\eqref{eq:crossing-momenta}
 &\(N=8\)&\(N=10\)&\(N=12\)&wrong \(G_0-G_1\), \(N=12\)\\
\midrule
first  &\(6.80\times10^{-5}\)&\(9.43\times10^{-6}\)&\(1.32\times10^{-6}\)&\(6.32\times10^{-2}\)\\
second &\(7.25\times10^{-5}\)&\(1.01\times10^{-5}\)&\(1.42\times10^{-6}\)&\(6.55\times10^{-2}\)\\
third  &\(7.92\times10^{-5}\)&\(1.12\times10^{-5}\)&\(1.58\times10^{-6}\)&\(7.34\times10^{-2}\)\\
\bottomrule
\end{tabular}
\end{center}
Here \(N\) is the maximum accumulated physical null level in the recursion.
The worst order-twelve point is the third case at \(z=0.8\).  At \(z=0.5\)
the three residuals are respectively \(7.11\times10^{-11}\),
\(2.61\times10^{-11}\), and \(2.02\times10^{-10}\); the largest residual
among the complex points is \(2.36\times10^{-8}\).  The maximum channel drift
from \(N=10\) to \(N=12\) is \(9.57\times10^{-6}\), showing that the edge
residual is dominated by recursion truncation.  At real \(z\), the largest
spurious imaginary part is \(7.46\times10^{-37}\).

As a negative control, the even and odd momentum integrals were also retained
separately.  At \(N=12\), the largest even-only and odd-only crossing
residuals are \(3.65\times10^{-2}\) and \(8.83\times10^{-1}\), whereas their
physical sum gives the table above.  Replacing that sum by the deliberately
wrong nonchiral combination \(G_0-G_1\) produces the last column.  This is a
test of the physical even/odd assembly, not a test of the chiral phase in
Eq.~\eqref{eq:odd-sign-cancels}.

Finally, for the first momentum set and
\(z=0.2,0.5,0.8,0.5+0.25i\), we varied the continuum integration at fixed
\(N=12\):
\begin{center}
\begin{tabular}{@{}lrr@{}}
\toprule
integration & maximum \(R\) & maximum channel change from previous row\\
\midrule
\(p\le4\), 24 points &\(1.319268\times10^{-6}\)&---\\
\(p\le5\), 32 points &\(1.319270\times10^{-6}\)&\(3.25\times10^{-9}\)\\
\(p\le6\), 40 points &\(1.319261\times10^{-6}\)&\(6.01\times10^{-11}\)\\
\bottomrule
\end{tabular}
\end{center}
At the last resolution, \(R(0.5)=1.40\times10^{-12}\) and
\(R(0.5+0.25i)=5.79\times10^{-12}\).  The integration drift is therefore
far below the order-twelve edge residual.  These are convergence diagnostics,
not certified error bounds for the infinite recursion.

The raw ledgers are
\begin{center}
\small
\path{Data Set/ns_liouville_crossing_extended_fresh.json},\\
\path{Data Set/ns_liouville_crossing_extended_p4_q24.json},\\
\path{Data Set/ns_liouville_crossing_extended_p6_q40.json}.
\end{center}
They support nonchiral crossing and its physical even/odd coefficients, but
the direct PBW comparison in the previous section remains what fixes the
human-note chiral sign \((-1)^{2n}\).

\section{Conclusion}

With the ungraded pairing, the auxiliary-fermion/SCA Gram matrix and the
double-Virasoro Gram matrix both factor without a parity insertion.  The
auxiliary identity block is one, so the two computations of the enlarged
sphere block give Eq.~\eqref{eq:final-B-formula}.  The only nontrivial
ordered-sphere sign is \((-1)^{2n}\), inherited from the human note's
third-slot odd tensor.  It passes the direct low-order computation; deleting
it produces the opposite odd block.  The extended crossing scan strongly
checks the physical nonchiral even/odd assembly and remains satisfied under
recursion and integration refinement, but it cannot distinguish the two
overall chiral odd-block sign conventions.

\end{document}
